\documentclass{jsarticle}
\usepackage{ascmac,amsmath,amssymb,amsthm,bm,physics,arydshln}
\newtheorem*{answer*}{解答}
\newtheorem*{another*}{別解}
\begin{document}

\begin{flushright}
6122111 三森誠真
\end{flushright}

\begin{itembox}[l]{\textbf{線形2.100}}
次の行列$A$と自然数$n$に対して, $A^n$を求めよ. \ 
$A=\left(\begin{array}{ccc}
5 & -1 & 1 \\
8 & -1 & 2 \\
-6 & 1 & -1
\end{array}\right)$.
\end{itembox}

\begin{answer*}
\end{answer*}
$$
\begin{aligned}
固有多項式を考えると, \quad|t E-A|=\left|\begin{array}{ccc}
t-5 & 1 & -1 \\
-8 & t+1 & -2 \\
6 & -1 & t+1
\end{array}\right| & =\left|\begin{array}{ccc}
t-5 & 1 & -1 \\
-(t-1)(t-3) & 0 & t-1 \\
t+1 & 0 & t
\end{array}\right| \\
 & =-\left|\begin{array}{cc}
-(t-1)(t-3) & t-1 \\
t+1 & t
\end{array}\right| \\
 \quad & =t(t-1)(t-3)+(t+1)(t-1) \\
& =(t-1)^3
\end{aligned}
$$



となり Cayley–Hamiltonの定理より$(A-E)^3=O$が成り立つので, 


$$
\begin{aligned}
A^n & = \{E+(A-E)\}^n \\
&=E+n(A-E)+\frac{n(n-1)}{2}(A-E)^2 \\
&=\frac{1}{2}\left(\begin{array}{lll}
2 & 0 & 0 \\
0 & 2 & 0 \\
0 & 0 & 2
\end{array}\right)+\frac{1}{2}\left(\begin{array}{ccc}
8 n & -2 n & 2 n \\
16 n & -4 n & 4 n \\
-12 n & 2 n & -4 n
\end{array}\right)+\frac{n(n-1)}{2}\left(\begin{array}{ccc}
2 & -1 & 0 \\
4 & -2 & 0 \\
-4 & 2 & 0
\end{array}\right) \\
&= \underline{\frac{1}{2}\left(\begin{array}{lll}
2n^2+6n+2 & -n^2-n & 2n \\
4n^2+12n & -2n^2-2n+2 & 4n \\
-4n^2-8n & 2n^2 & -4n+2
\end{array}\right)}.
\end{aligned}
$$

\bigskip
\bigskip

\begin{another*}
\end{another*}
$A$の固有多項式は$(t-1)^3$であり, \ 
$A-E=\left(\begin{array}{ccc}
4 & -1 & 1 \\
8 & -2 & 2 \\
-6 & 1 & -2
\end{array}\right) \neq O$, \ \ 
$(A-E)^2=\left(\begin{array}{ccc}
2 & -1 & 0 \\
4 & -2 & 0 \\
-4 & 2 & 0
\end{array}\right) \neq O$

なので最小多項式は$(t-1)^3$となり$A$は対角化不可能.

そこで, $A$のJordan標準形を求めていく.

$$
\begin{gathered}
\left(\begin{array}{c:c}
    & x_1 \\
A-E & x_2 \\
    & x_3
\end{array}\right) \rightarrow
\left(\begin{array}{ccc:c}
4 & -1 & 1 & x_1 \\
8 & -2 & 2 & x_2 \\
-6 & 1 & -2 & x_3
\end{array}\right) \rightarrow
\left(\begin{array}{ccc:c}
1 & 0 & \frac{1}{2} & -\frac{1}{2} x_1-\frac{1}{2} x_3 \\
0 & 1 & 1 & -3 x_1-2 x_3 \\
0 & 0 & 0 & -2 x_1+x_2
\end{array}\right) \\
\end{gathered}
$$

となるので, \ 
$(A-E) \bm{p}_1 = \bm{0}$を解くと, 
$\bm{p}_1=
\left(\begin{array}{c}
1 \\
2 \\
-2
\end{array}\right).$ \\
次に, $(A-E) \bm{p}_2=\bm{p}_1$を解くと,$\bm{p}_2=\left(\begin{array}{c}
-\frac{1}{2} c+\frac{1}{2} \\
-c+1 \\
c
\end{array}\right)
(c \in \mathbb{R})
$となり, 
$c=1$を代入すると, 
$\bm{p}_2=\left(\begin{array}{c}
0 \\
0 \\
1
\end{array}\right)$ \\
となる. \\
最後に,$(A-E) \bm{p}_3=\bm{p}_2$を解くと,$\bm{p}_3=\left(\begin{array}{c}
-\frac{1}{2} d-\frac{1}{2} \\
-d-2 \\
d
\end{array}\right)(d \in \mathbb{R})$となり, 
$d=-1$を代入すると, 
$\bm{p}_3=\left(\begin{array}{c}
0 \\
-1 \\
-1
\end{array}\right)$
となる.


\( \begin{aligned} & P:=\left(\begin{array}{ccc}\bm{p}_1 & \bm{p}_2 & \bm{p}_3\end{array}\right)=\left(\begin{array}{ccc}1 & 0 & 0 \\ 2 & 0 & -1 \\ -2 & 1 & -1\end{array}\right) \text{とおくと, } \\ & P^{-1} A P=\left(\begin{array}{ccc}1 & 0 & 0 \\ 4 & -1 & 1 \\ 2 & -1 & 0\end{array}\right)\left(\begin{array}{ccc}5 & -1 & 1 \\ 8 & -1 & 2 \\ -6 & 1 & -1\end{array}\right)\left(\begin{array}{ccc}1 & 0 & 0 \\ 2 & 0 & -1 \\ -2 & 1 & -1\end{array}\right) =\left(\begin{array}{ccc}5 & -1 & 1 \\ 6 & -2 & 1 \\ 2 & -1 & 0\end{array}\right)\left(\begin{array}{ccc}1 & 0 & 0 \\ 2 & 0 & -1 \\ -2 & 1 & -1\end{array}\right) \\ & \hspace{31zw}=\left(\begin{array}{ccc}1 & 1 & 0 \\ 0 & 1 & 1 \\ 0 & 0 & 1\end{array}\right) \\ &\hspace{31zw}=E+\left(\begin{array}{ccc}0 & 1 & 0 \\ 0 & 0 & 1 \\ 0 & 0 & 0\end{array}\right) \\
&\text{となり, 両辺を}n\text{乗すると} \\ 
&\left(P^{-1} A P\right)^n=E^n+n\left(\begin{array}{ccc}0 & 1 & 0 \\ 0 & 0 & 1 \\ 0 & 0 & 0\end{array}\right) + \displaystyle\frac{n(n-1)}{2}\left(\begin{array}{ccc}0 & 1 & 0 \\ 0 & 0 & 1 \\ 0 & 0 & 0\end{array}\right)^2 \hspace{5zw}\left(\because \left(\begin{array}{ccc}0 & 1 & 0 \\ 0 & 0 & 1 \\ 0 & 0 & 0\end{array}\right)^3 = O \right)\\ &\hspace{5.5zw}= E+\left(\begin{array}{ccc}0 & n & 0 \\ 0 & 0 & n \\ 0 & 0 & 0\end{array}\right)+\frac{1}{2}\left(\begin{array}{ccc}0 & 0 & n^2-n \\ 0 & 0 & 0 \\ 0 & 0 & 0\end{array}\right) \\ &\hspace{5.5zw}=\frac{1}{2}\left(\begin{array}{ccc}2 & 2 n & n^2-n \\ 0 & 2 & 2 n \\ 0 & 0 & 2\end{array}\right). & \end{aligned} \) 

\bigskip

よって, $A^n = P \ (P^{-1}AP)^n \ P^{-1}$ \\
\hspace{6.5zw}= \( \displaystyle\frac{1}{2}\left(\begin{array}{ccc}2 & 2 n & n^2-n \\ 4 & 4 n & 2 n^2-2 n-2 \\ -4 & 2-4 n & -2 n^2+4 n-2\end{array}\right)\left(\begin{array}{ccc}1 & 0 & 0 \\ 4 & -1 & 1 \\ 2 & -1 & 0\end{array}\right)\) \\

$\hspace{5.5zw}=\displaystyle \underline{\frac{1}{2} \begin{pmatrix}
2n^2+6n+2 & -n^2-n & 2n \\
4n^2+12n & -2n^2-2n+2 & 4n \\
-4n^2-8n & 2n^2 & -4n+2 \end{pmatrix}}.$

\bigskip
\bigskip

\begin{another*}
\end{another*}

$W(1;A) = \left< \begin{pmatrix} 1 \\ 2 \\ -2 \end{pmatrix} \right>$ なので, ある正規直交基底$\left\{\displaystyle \frac{1}{3} \begin{pmatrix} 1 \\ 2 \\ -2 \end{pmatrix}, \ 
        \displaystyle \frac{1}{\sqrt{2}} \begin{pmatrix} 0 \\ 1 \\ 1 \end{pmatrix}, \ 
        \displaystyle \frac{1}{3\sqrt{2}} \begin{pmatrix} -4 \\ 1 \\ -1 \end{pmatrix} \right\}$
を考えて, 

$P:= \displaystyle\frac{1}{3\sqrt{2}} 
\begin{pmatrix}  \sqrt{2} & 0 & -4 \\ 2\sqrt{2} & 3 & 1 \\ -2\sqrt{2} & 3 & -1 \end{pmatrix}$ とおくと, 

$P^{-1}AP = {}^t PAP = \displaystyle\frac{1}{18} 
\begin{pmatrix}  \sqrt{2} & 2\sqrt{2} & -2\sqrt{2} \\ 0 & 3 & 3 \\ -4 & 1 & -1 \end{pmatrix}
\left(\begin{array}{ccc}
5 & -1 & 1 \\
8 & -1 & 2 \\
-6 & 1 & -1
\end{array}\right)
\begin{pmatrix}  \sqrt{2} & 0 & -4 \\ 2\sqrt{2} & 3 & 1 \\ -2\sqrt{2} & 3 & -1 \end{pmatrix} \\
\hspace{9zw} = \displaystyle\frac{1}{18} 
\begin{pmatrix}  33\sqrt{2} & -5\sqrt{2} & 7\sqrt{2} \\ 6 & 0 & 3 \\ -6 & 2 & -1 \end{pmatrix}
\begin{pmatrix}  \sqrt{2} & 0 & -4 \\ 2\sqrt{2} & 3 & 1 \\ -2\sqrt{2} & 3 & -1 \end{pmatrix} \\
\hspace{9zw} = \displaystyle\frac{1}{6} \begin{pmatrix} 6 & 2\sqrt{2} & -48\sqrt{2} \\ 0 & 3 & -9 \\ 0 & 1 & 9 \end{pmatrix}$.

次に, $B:= \begin{pmatrix} 3 & -9 \\ 1 & 9 \end{pmatrix}$ とおくと, $|tE-B| = \begin{vmatrix} t-3 & 9 \\ -1 & t-9 \end{vmatrix} = (t-6)^2$ より固有値が$6$であり, 

$6E-B = \begin{pmatrix} 3 & 9 \\ -1 & -3 \end{pmatrix} $なので, $W(6;A) = \left< \begin{pmatrix} 3 \\ -1 \end{pmatrix} \right>$. よって, 正規直交基底$\left\{\displaystyle\frac{1}{\sqrt{10}}\begin{pmatrix} 3 \\ -1 \end{pmatrix} , \displaystyle\frac{1}{\sqrt{10}}\begin{pmatrix} 1 \\ 3 \end{pmatrix} \right\}$ 

を考えて$Q := \displaystyle\frac{1}{\sqrt{10}}\begin{pmatrix} 3 & 1 \\ -1 & 3 \end{pmatrix}$ とおくと, 
$Q^{-1}BQ = \displaystyle\frac{1}{10} \begin{pmatrix} 3 & -1 \\ 1 & 3 \end{pmatrix}$
$\begin{pmatrix} 3 & -9 \\ 1 & 9 \end{pmatrix} \begin{pmatrix} 3 & 1 \\ -1 & 3 \end{pmatrix}$

$\hspace{23zw} = \displaystyle\frac{1}{10} \begin{pmatrix} 8 & -36 \\ 6 & 18 \end{pmatrix} \begin{pmatrix} 3 & 1 \\ -1 & 3 \end{pmatrix}$

$\hspace{23zw} = \begin{pmatrix} 6 & -10 \\ 0 & 6 \end{pmatrix}$.

$R:= P \left( 
\begin{array}{c:c}
1 & \bm{0} \\ \hdashline
 &  \\ 
\bm{0} & \ Q  \\
 &  \\
\end{array}
\right) = \displaystyle \frac{1}{3\sqrt{5}}\begin{pmatrix} \sqrt{5} & 2 & -6 \\ 2\sqrt{5} & 4 & 3 \\ -2\sqrt{5} & 5 & 0 \end{pmatrix}$ とおくと, 

$R^{-1}AR = \left( \begin{array}{c:c} 1 & \bm{0} \\ \hdashline &  \\ 
\bm{0} & \ Q^{-1}  \\
 &  \\
\end{array}
\right) P^{-1}AP
\left( 
\begin{array}{c:c}
1 & \bm{0} \\ \hdashline
 &  \\ 
\bm{0} & \ Q  \\
 &  \\
\end{array}
\right)$
= $\displaystyle\frac{1}{6} \left( \begin{array}{c:c} 1 & \bm{0} \\ \hdashline &  \\ 
\bm{0} & \ Q^{-1}  \\
 &  \\
\end{array}
\right) 
\left(
\begin{array}{c:cc} 6 & 2\sqrt{2} & -48\sqrt{2} \\ \hdashline & & \\ \bm{0} & \ B & \\ & &   \end{array} 
\right)
\left( 
\begin{array}{c:c}
1 & \bm{0} \\ \hdashline
 &  \\ 
\bm{0} & \ Q  \\
 &  \\
\end{array}
\right)$

$\hspace{22zw} = \displaystyle\frac{1}{6} \left( \begin{array}{c:cc} 6 & 2\sqrt{2} & -48\sqrt{2} \\ \hdashline &  \\ \bm{0} & \ Q^{-1}B  \\
 &  \\
\end{array} \right)
\left( 
\begin{array}{c:c}
1 & \bm{0} \\ \hdashline
 &  \\ 
\bm{0} & \ Q  \\
 &  \\
\end{array}
\right)$

$\hspace{22zw} = \begin{pmatrix} 1 & \displaystyle\frac{9}{\sqrt{5}} & -\displaystyle\frac{71}{3\sqrt{5}} \\ 0 & 1 & -\displaystyle\frac{5}{3} \\ 0 & 0 & 1 \end{pmatrix}$

$\hspace{22zw} = E + \begin{pmatrix} 0 & \displaystyle\frac{9}{\sqrt{5}} & -\displaystyle\frac{71}{3\sqrt{5}} \\ 0 & 0 & -\displaystyle\frac{5}{3} \\ 0 & 0 & 0 \end{pmatrix}$.

\bigskip

$(R^{-1}AR)^n = E^n + nE^{n-1} \begin{pmatrix} 0 & \displaystyle\frac{9}{\sqrt{5}} & -\displaystyle\frac{71}{3\sqrt{5}} \\ 0 & 0 & -\displaystyle\frac{5}{3} \\ 0 & 0 & 0 \end{pmatrix} + \displaystyle \frac{n(n-1)}{2} E^{n-2} \begin{pmatrix} 0 & \displaystyle\frac{9}{\sqrt{5}} & -\displaystyle\frac{71}{3\sqrt{5}} \\ 0 & 0 & -\displaystyle\frac{5}{3} \\ 0 & 0 & 0 \end{pmatrix}^2$

= $\displaystyle\frac{1}{6\sqrt{5}} \begin{pmatrix} 6\sqrt{5} & 0 & 0 \\ 0 & 6\sqrt{5} & 0 \\ 0 & 0 & 6\sqrt{5} \end{pmatrix} + \displaystyle\frac{1}{6\sqrt{5}} \begin{pmatrix} 0 & 54n & -142n \\ 0 & 0 & -10\sqrt{5}n \\ 0 & 0 & 0 \end{pmatrix} - \displaystyle\frac{1}{6\sqrt{5}} \begin{pmatrix} 0 & 0 & 45n(n-1) \\ 0 & 0 & 0 \\ 0 & 0 & 0 \end{pmatrix}$

= $\displaystyle\frac{1}{6\sqrt{5}} \begin{pmatrix} 6\sqrt{5} & 54n & -45n^2-97n \\ 0 & 6\sqrt{5} & -10\sqrt{5} \\ 0 & 0 & 6\sqrt{5} \end{pmatrix}$



\end{document}
